% HD-I24. GPT-6 Astra Pro /root/imprimitive; inherited effort, not numerically exposed.
% Author draft, 21 November 2026. Exact correlated quotient, not a product replacement.
\section{HD-I24: the correlated tensor-two endpoint}
\label{sec:correlated-tensor-two-rank-four}

Let \(K\) be an imaginary quadratic field. Consider an actual finite
subgroup
\[
 H\subset S_s\times GL(M),\qquad \dim_K M=2,\qquad
 A_s\times1\subset H,\qquad s=r+1,\quad r\ge11.             \tag{I24.1}
\]
On \(W=\{(x_1,\ldots,x_s)\in K^s:\sum x_i=0\}\), use the
literal permutation action; the given coefficient representation is
\(V=W\otimes_K M\), with action \(P_\pi\otimes B\).
Assume it is faithful and absolutely irreducible, has exact character
field \(K\), and has no nonidentity element of moving rank one, two,
or three. A connected full-monodromy sphere cover with this coefficient
representation and a prescribed primitive factor of signature \((3,3)\)
has period rank at most eight. The conclusion concerns its original
corner and restricted polarization. In particular this endpoint cannot
produce a missing period-rank-nine natural source.

The projection \(\Gamma\subset GL(M)\) is the actual image of \(H\).
It is finite and absolutely irreducible: an invariant proper complex
subspace of \(M\) would give one in \(V\). Its trace field is \(K\):
every trace on \(V\) is the product of the integral augmentation trace
and a trace on \(M\), so
\(K=\mathbf Q(\chi_V(H))\subset\mathbf Q(\chi_M(\Gamma))\subset K\).
We retain \(H\) itself throughout; it can be the subgroup defined by a
common sign character, and need not be a direct product.

\paragraph{1. The exact inertia budget.}
Let \(b\) be the number of nontrivial coefficient branch values, and
\(c_j\) their moving ranks. Irreducibility kills the invariant and
coinvariant terms in the cellular Euler formula on the punctured
sphere. Thus the first parabolic cohomology of dimension six gives
\[
                 \sum_jc_j=4r+6.                           \tag{I24.2}
\]
Period rank nine would require \(b\ge12\), because the marked sphere
has dimension \(b-3\). We prove that this requirement forces a
factorization through a period space of dimension at most eight.

Let \(k\) count the branch entries whose image in \(PGL(M)\) is
nonidentity. There are at least three. Otherwise the product-one
relation makes their projective images generate a cyclic group;
its finite linear inverse image consists of powers of one diagonalizable
matrix and scalar matrices, so it preserves a complex line.
For nonscalar \(B\), its two distinct eigenvalues show directly that
\(\dim\ker(P_\pi\otimes B-I)\le r\): the two matching eigenspaces
of \(P_\pi\) are disjoint. Such an entry has \(c_j\ge r\).
Every other nonidentity coefficient entry has \(c_j\ge4\).
Consequently
\[
        4r+6\ge kr+4(b-k)\ge k(r-4)+48.                    \tag{I24.3}
\]
Since \(r\ge11\), \(k\ge4\) is impossible. Therefore
\(k=3\), and the same inequality also gives \(r\ge30\).

All remaining images in \(\Gamma\) are the identity, rather than
nonidentity scalar matrices. Indeed, for a scalar \(aI\ne I\), the
\(a^{-1}\)-eigenspace of a permutation on \(K^s\) has dimension at
most \(\lfloor s/2\rfloor\): every cycle supporting a nonidentity
eigenvalue has length at least two and contributes at most one.
The same bound applies to its augmentation subspace. Thus the
coefficient moving rank is at least
\(2r-2\lfloor s/2\rfloor\ge r-1\). Together with the three
nonscalar entries and at least eight further entries, this gives
\(\sum c_j\ge3r+(r-1)+32=4r+31\), contrary to (I24.2).

The projected \(\Gamma\)-tuple therefore has exactly three
nonidentity entries, all nonscalar. Its actual rank-two parabolic
cohomology on the sphere has dimension
\[
 d_0=-4+\sum_{j=1}^3\operatorname{rank}(\gamma_j-I)
                         \in\{0,1,2\}.                    \tag{I24.4}
\]
Nonnegativity follows from this being a cohomology dimension, not
from a separate tuple assumption.

\paragraph{2. The case \(d_0=2\) is impossible.}
Here all three \(\gamma_j\) are fixed-point-free and nonscalar.
Put \(f_\ell=\dim\ker(\gamma_j^\ell-I)\). Then \(f_1=0\)
and \(f_2\le1\): equality \(f_2=2\) would make a fixed-point-free
order-two matrix the scalar \(-I\). For \(\ell\ge3\),
\(f_\ell\le2\). Hence \(f_\ell\le2\ell/3\) for every
positive \(\ell\). On a permutation cycle of length \(\ell\),
the fixed-vector equations identify the fixed space of
\(P_\pi\otimes\gamma_j\) with \(\ker(\gamma_j^\ell-I)\).
Removing the constant permutation summand subtracts \(f_1=0\).
It follows that
\[
 \dim\ker(P_\pi|_W\otimes\gamma_j-I)\le\frac{2s}{3},
 \qquad c_j\ge\frac{4r-2}{3}.                             \tag{I24.5}
\]
The three such entries and at least nine others now give
\(\sum c_j\ge4r-2+36=4r+34\), again contradicting (I24.2).
Thus \(d_0\le1\).

\paragraph{3. The literal quotient curves, including the sign correlation.}
Write \(\widetilde C\) for the given full-\(H\) curve. Set
\[
 K_0=\ker(H\longrightarrow\Gamma),\quad
 J=\{(\pi,B)\in H:\pi(1)=1\},\quad L=J\cap K_0.
                                                               \tag{I24.6}
\]
The group \(K_0\) is either \(A_s\times1\) or \(S_s\times1\),
and \(L\) is its point stabilizer. The map \(J\to\Gamma\) is
onto: multiply any lift by an element of \(A_s\times1\) to make
its permutation fix 1. Thus \(L\lhd J\), \(J/L=\Gamma\),
\(JK_0=H\), and \([K_0:L]=s\). These are statements about
the specified subgroup, without choosing a section of \(H\to\Gamma\).
Take the actual curves and finite maps
\[
 X=\widetilde C/L\longrightarrow Y=\widetilde C/J,
 \qquad Z=\widetilde C/K_0\longrightarrow\mathbf P^1,
 \qquad f:X\longrightarrow Z.                              \tag{I24.7}
\]
Both displayed covers are full-\(\Gamma\) covers. The degree of
\(f\) is \(s\), and \(f\) is \(\Gamma\)-equivariant. Indeed
the action on \(X\) is through \(J/L\), the action on \(Z\)
through \(H/K_0\), and these quotient maps agree. Their function
fields show that \(X\) is the normalization of
\(Y\times_{\mathbf P^1} Z\): the relevant subgroup intersection
is \(J\cap K_0=L\), and their product is \(H\).

The actual induction isomorphism is
\[
 \operatorname{Ind}_J^H(M)
    \simeq K^s\otimes M=(K\mathbf1\otimes M)\oplus(W\otimes M),
                                                               \tag{I24.8}
\]
where the sheet \(i\) carries \(M\) and an element
\((\pi,B)\) sends \(e_i\otimes u\) to \(e_{\pi(i)}\otimes Bu\).
The summand maps are the diagonal inclusion and the average
\(s^{-1}\sum_i\), with their explicit complementary projection.
At a branch cycle, the fixed-vector equations on a cycle of length
\(\ell\) give \(\ker(\gamma^\ell-I)\) on both sides. Thus
(I24.8) identifies the middle extensions and their first parabolic
cohomology, including the ramified stalks. In particular the full
rank-two system on \(Y\) has parabolic dimension
\[
                              6+d_0\le7.                   \tag{I24.9}
\]

\paragraph{4. The actual corner and its polarization.}
Choose a positive invariant \(K\)-Hermitian form \(h_M\) on
\(M\). The invariant form on \(V\), restricted from the original
coefficient system, is a positive rational scalar times
\(h_W\otimes h_M\), after rescaling \(h_M\) if needed.
This follows directly from the commutant of the absolutely
irreducible rational \(A_s\)-module \(W\): in a matrix basis,
an \(A_s\)-invariant form is \(h_W\) tensored with a form on
\(M\); invariance under \(H\) is precisely invariance of that
second form under \(\Gamma\). Here \(h_W\) is the restriction
of \(\sum x_i\bar y_i\).

For \(u\in M(K)\setminus\{0\}\), let \(e_u\in\mathbf Q[\Gamma]\)
be its rational orthogonal rank-one corner (including the conjugate
complex representation). Its action on \(JX\) commutes with
\(f^*f_*\). Define the actual rational Jacobian projector
\[
                \pi=e_u\left(1-\frac{f^*f_*}{s}\right).
                                                               \tag{I24.10}
\]
Since \(f_*f^*=[s]\), this is a Rosati-self-adjoint idempotent.
For complete group-algebra precision, inflate \(e_u\) to
\(\widetilde e_u\in\mathbf Q[J]\) by replacing each
\(\gamma\) by \(\operatorname{av}_L j\), for any lift
\(j\in J\), where \(\operatorname{av}_L=|L|^{-1}\sum_{l\in L}l\).
This is independent of the lift. With \(\alpha:\widetilde C\to X\),
pullback and norm give the literal operator identity, with its
specified lift to the group algebra,
\[
 \frac{\alpha^*\pi\alpha_*}{|L|}
                =(1-\operatorname{av}_{K_0})\widetilde e_u
                     \quad\bigl(\text{right side in }\mathbf Q[H]\bigr).
                                                               \tag{I24.11}
\]
For the second term, use
\((f\alpha)^*(f\alpha)_*=
 |K_0|\operatorname{av}_{K_0}\), and
\(|K_0|=s|L|\); equivariance moves \(e_u\) through \(f_*\).
This verifies the identity without a product decomposition of \(H\).

By (I24.8), the inflated \(J\)-corner is supported on precisely
the two indicated \(H\)-representations, with their conjugates.
On \(W\otimes M\), \(\operatorname{av}_L\) is the orthogonal
projection to
\(Kw_0\otimes M\), where \(w_0=(s-1,-1,\ldots,-1)\).
The subtraction in (I24.11) removes exactly the constant summand.
Thus its right side is the original rational-block primitive
corner for \(w_0\otimes u\), not an additional factor.
Clearing the denominators of \(\pi\), \(\alpha^*\) restricts
to a surjective isogeny from its connected image in \(JX\)
onto that corner of \(J\widetilde C\), with exact polarization
multiplier \(|L|\):
\[
                 \alpha_*\alpha^*=[|L|],\qquad
                 (\alpha^*)^*\theta_{J\widetilde C}
                                      =|L|\theta_{JX}.       \tag{I24.12}
\]
An initially prescribed coefficient line \(Kv\) is changed to
\(K(w_0\otimes u)\) by the actual rational matrix-unit element
of HD-G29, (G29.15a): for \(T(x)=w_0\otimes u\,h(v,x)/h(v,v)\),
\[
 b=\frac{2\dim_K V}{|H|}\sum_{h\in H}
       \operatorname{Re}\operatorname{tr}(T\rho(h^{-1}))h,
 \qquad b^\dagger b=
 \frac{h(w_0\otimes u,w_0\otimes u)}{h(v,v)}e_v.            \tag{I24.13}
\]
An integer-cleared \(b\) gives a polarized isogeny with the
displayed positive rational multiplier times an integer square.
All maps retain the actual \(K\)-action. Their projector matrices
are flat after fixing the finite cover markings; hence the
original period differential factors through that of the full
rank-two cover system on \(Y\).

\paragraph{5. The period bound for that full system.}
Put \(a=6+d_0\le7\), let \(g=g(Y)\), and omit points with
identity coefficient inertia. The same Euler computation gives
\[
                       a=2(2g-2)+\sum_{j=1}^{b_Y}c_j,
                       \qquad c_j\ge1.                    \tag{I24.14}
\]
For \(g=0\), \(b_Y\le a+4\le11\), so the marked-base
dimension is at most eight. For \(g=1\) it is at most seven;
for \(g=2\), \(b_Y\le3\), giving at most six. Genus at
least three is impossible. Finite-monodromy covers have discrete
choices over a fixed pointed base. The entire quotient period map,
and therefore its flat projector factor (I24.10), has rank at
most eight. This proves the asserted source exclusion.
